% GENERATED FILE. Edit canonical lessons, then run npm run generate:books.
% canonical-source-hash: fnv1a64:fb116499c8a552f2
% canonical-lessons: GE-C236-pflicht, GE-C236-qualitaet, GE-C236-rede, GE-C236-reihe, GE-C236-religion

\chapter{A duty, Quality, A speech, A row, Religion}
\label{ch:ge-a-duty-quality-a-speech-a-row-religion}
\hlchaptermodality{236}

\begin{chapteropening}
\textbf{By the end of this chapter:} \emph{I can say a duty, quality, a speech, a row and religion.}

Complete the last lesson of chapter 236: I can say a duty, quality, a speech, a row and religion.
\end{chapteropening}

\section[die Pflicht]{die Pflicht — a duty}

\label{lesson:GE-C236-pflicht}



\begin{quote}
Before the new one: say the German for order, then the German for a couple.
\end{quote}

\subsection*{You'll want to know: die Pflicht}
\textbf{die Pflicht} — \textquotedblleft{}a duty\textquotedblright{}.

\begin{cousinweb}
\textbf{Pflicht} $\leftarrow$ Old High German \emph{pfliht} $\leftarrow$ Proto-Germanic \emph{*plihtiz} — inherited, the same word as English \emph{plight}, as in \textquotedblleft{}plight one's troth\textquotedblright{}. (EtymDB 2.1)
\end{cousinweb}

\begin{grammarlens}[title={What you've built}]
One more word for this chapter.
\end{grammarlens}

\subsection*{Your turn}
\begin{itemize}
\raggedright
  \item \emph{Say it:} \emph{die Pflicht}
  \item \emph{Say it:} \emph{die Pflicht}, once more
  \item \emph{Say it:} say \emph{die Pflicht}
  \item \emph{Recall:} say the German for order, then the German for a couple, then say \emph{die Pflicht} again
\end{itemize}

\begin{tcolorbox}[breakable,colback=teal!4,colframe=teal!35!black,title={Before you move on}]
What does die Pflicht mean? (\textquotedblleft{}A duty\textquotedblright{}.) Say it once more.
\end{tcolorbox}
\section[die Qualität]{die Qualität — quality}

\label{lesson:GE-C236-qualitaet}



\begin{quote}
Before the new one: say the German for a couple, then the German for a duty.
\end{quote}

\subsection*{You'll want to know: die Qualität}
\textbf{die Qualität} — \textquotedblleft{}quality\textquotedblright{}.

\begin{cousinweb}
\textbf{Qualität} $\leftarrow$ Latin \emph{quālitās}, \textquotedblleft{}what kind something is\textquotedblright{}, from \emph{quālis}, \textquotedblleft{}of what kind\textquotedblright{} (EtymDB 2.1) — borrowed, like English \emph{quality}.
\end{cousinweb}

\begin{grammarlens}[title={What you've built}]
One more word for this chapter.
\end{grammarlens}

\subsection*{Your turn}
\begin{itemize}
\raggedright
  \item \emph{Say it:} \emph{die Qualität}
  \item \emph{Say it:} \emph{die Qualität}, once more
  \item \emph{Say it:} say \emph{die Qualität}
  \item \emph{Recall:} say the German for a couple, then the German for a duty, then say \emph{die Qualität} again
\end{itemize}

\begin{tcolorbox}[breakable,colback=teal!4,colframe=teal!35!black,title={Before you move on}]
What does die Qualität mean? (\textquotedblleft{}Quality\textquotedblright{}.) Say it once more.
\end{tcolorbox}
\section[die Rede]{die Rede — a speech}

\label{lesson:GE-C236-rede}



\begin{quote}
Before the new one: say the German for a duty, then the German for quality.
\end{quote}

\subsection*{You'll want to know: die Rede}
\textbf{die Rede} — \textquotedblleft{}a speech\textquotedblright{}.

\begin{cousinweb}
\textbf{Rede} $\leftarrow$ Middle High German \emph{rede} $\leftarrow$ Old High German \emph{redia}, \textquotedblleft{}speech, account\textquotedblright{} (EtymDB 2.1) — inherited.
\end{cousinweb}

\begin{grammarlens}[title={What you've built}]
One more word for this chapter.
\end{grammarlens}

\subsection*{Your turn}
\begin{itemize}
\raggedright
  \item \emph{Say it:} \emph{die Rede}
  \item \emph{Say it:} \emph{die Rede}, once more
  \item \emph{Say it:} say \emph{die Rede}
  \item \emph{Recall:} say the German for a duty, then the German for quality, then say \emph{die Rede} again
\end{itemize}

\begin{tcolorbox}[breakable,colback=teal!4,colframe=teal!35!black,title={Before you move on}]
What does die Rede mean? (\textquotedblleft{}A speech\textquotedblright{}.) Say it once more.
\end{tcolorbox}
\section[die Reihe]{die Reihe — a row, a turn}

\label{lesson:GE-C236-reihe}



\begin{quote}
Before the new one: say the German for quality, then the German for a speech.
\end{quote}

\subsection*{You'll want to know: die Reihe}
\textbf{die Reihe} — \textquotedblleft{}a row, a turn\textquotedblright{}.

\begin{cousinweb}
\textbf{Reihe} $\leftarrow$ Middle High German \emph{rīhe}, \textquotedblleft{}row, line\textquotedblright{} (EtymDB 2.1) — inherited. Your turn is your place in the row.
\end{cousinweb}

\begin{grammarlens}[title={What you've built}]
One more word for this chapter.
\end{grammarlens}

\subsection*{Your turn}
\begin{itemize}
\raggedright
  \item \emph{Say it:} \emph{die Reihe}
  \item \emph{Say it:} \emph{die Reihe}, once more
  \item \emph{Say it:} say \emph{die Reihe}
  \item \emph{Recall:} say the German for quality, then the German for a speech, then say \emph{die Reihe} again
\end{itemize}

\begin{tcolorbox}[breakable,colback=teal!4,colframe=teal!35!black,title={Before you move on}]
What does die Reihe mean? (\textquotedblleft{}A row, a turn\textquotedblright{}.) Say it once more.
\end{tcolorbox}
\section[die Religion]{die Religion — religion}

\label{lesson:GE-C236-religion}



\begin{quote}
Before the new one: say the German for a speech, then the German for a row.
\end{quote}

\subsection*{You'll want to know: die Religion}
\textbf{die Religion} — \textquotedblleft{}religion\textquotedblright{}.

\begin{cousinweb}
\textbf{Religion} $\leftarrow$ Latin \emph{religiō}, \textquotedblleft{}scruple, reverence\textquotedblright{} (EtymDB 2.1) — borrowed, like English \emph{religion}.
\end{cousinweb}

\begin{grammarlens}[title={What you've built}]
One more word for this chapter.
\end{grammarlens}

\subsection*{Your turn}
\begin{itemize}
\raggedright
  \item \emph{Say it:} \emph{die Religion}
  \item \emph{Say it:} \emph{die Religion}, once more
  \item \emph{Say it:} say \emph{die Religion}
  \item \emph{Recall:} say the German for a speech, then the German for a row, then say \emph{die Religion} again
\end{itemize}

\begin{tcolorbox}[breakable,colback=teal!4,colframe=teal!35!black,title={Before you move on}]
What does die Religion mean? (\textquotedblleft{}Religion\textquotedblright{}.) Say it once more.
\end{tcolorbox}
